\begin{solapartat}
\punts{0,5} A partir del nivell d'intensitat sonora obtenim la intensitat de l'ona: $\beta = 10\log(\frac{I}{I_0})$, $100 = 10\log(\frac{I}{10^{-12}})$, $I = 10^{10}\cdot 10^{-12} = 10^{-2}\ \mathrm{W/m^2}$.

\punts{0,5} La potència emesa pel petard suposant que el so es reparteix en una superfície esfèrica: $P = I\cdot 4\pi r^2 = 10^{-2}4\pi\cdot 50^2 = 314,16\ \mathrm{W}$.

\punts{0,25} L'energia sonora alliberada: $E = P\cdot t = 314,16\cdot 0,03 = 9,42\ \mathrm{J}$.
\end{solapartat}

\begin{solapartat}
\punts{0,25} A partir del nivell d'intensitat sonora obtenim la intensitat de l'ona: $\beta = 10\log(\frac{I}{I_0})$, $90 = 10\log(\frac{I}{10^{-12}})$, $I = 10^{9}\cdot 10^{-12} = 10^{-3}\ \mathrm{W/m^2}$.

\textit{Càlcul de l'alçada:}

\punts{0,25} Els coets són iguals i, per tant, la potència emesa també. Es reparteix en una superfície esfèrica i obtenim la distància a la qual es trobava el petard utilitzant la mateixa relació:
\[ P = I\cdot 4\pi r^2,\quad r^2 = \frac{P}{4\pi\cdot I} = \frac{314,16}{4\pi\cdot 10^{-3}} = 25000\ \mathrm{m^2},\quad r = 158,11\ \mathrm{m} \]

\punts{0,25} Utilitzant el teorema de Pitàgores obtenim l'alçada a la qual esclata el coet:
\[ h = \sqrt{158,11^2 - 50^2} = 150\ \mathrm{m} \]

\textit{O alternativament:}

\punts{0,5} Els coets són iguals i, per tant, la potència emesa també. Es reparteix en una superfície esfèrica i, per tant, entre les intensitats i les distàncies hi ha la relació següent:
\[ \frac{I_2}{I_1} = \frac{r_1^2}{r_2^2} \text{ per tant } \frac{10^{-3}\ \mathrm{W/m^2}}{10^{-2}\ \mathrm{W/m^2}} = 0,1 = \frac{50^2}{h^2 + 50^2} \text{ i per tant } h^2 + 50^2 = 10\cdot 50^2, \]
$h^2 = 9\cdot 50^2$, $h = 3\cdot 50 = 150\ \mathrm{m}$

\punts{0,5} La intensitat sonora generada per dos coets pot ser calculada sabent que la potencia i la intensitat del so es doblaran. Per tant: $\beta = 10\log\left(\frac{2\cdot 10^{-3}}{10^{-12}}\right) = 93\ \mathrm{dB}$.
\end{solapartat}
